\documentclass[10pt,a4paper]{amsart}
\usepackage[margin=19mm]{geometry}
\usepackage{amsmath,amssymb,mathtools}
\usepackage[scr=rsfs,cal=euler]{mathalfa}
\usepackage{xfrac}
\usepackage[unicode,hidelinks,bookmarks=false]{hyperref}
\newcommand{\ie}{i.e.\ }
\newcommand{\cf}{cf.\ }
\newcommand{\resp}{respectively\ }
\newcommand{\Subsection}{\subsection}
\providecommand{\nobreakdash}{\nobreak\hbox{-}}
\setlength{\emergencystretch}{2em}
\multlinegap0pt
\usepackage{amssymb}
\usepackage{mathtools}
\usepackage{dsfont}
\usepackage[shortlabels]{enumitem}
\usepackage{float}
\newtheorem{theorem}{Theorem}
\title{Discrete Hyperbolic Secant Distributions}
\author{Leonard Pleschberger}
\date{Source: arXiv:2609.11293v1 (CC BY 4.0)}
\begin{document}
\maketitle

\noindent\emph{Source and licence.} Discrete Hyperbolic Secant Distributions. Version arXiv:2609.11293v1 (CC BY 4.0), distributed by arXiv under the Creative Commons Attribution 4.0 International licence, http://creativecommons.org/licenses/by/4.0/, as recorded in the arXiv licence metadata for this exact version and shown on the abstract page https://arxiv.org/abs/2609.11293v1. Full licence notice: \texttt{licenses/arxiv-2609.11293-CC-BY-4.0.txt}.\par\smallskip

\noindent\textit{Editorial scope.} Counted unit: the article's theorem theorem (source0.tex lines 348--358). The article's complete original proof of it is delivered with it (source0.tex lines 360--376, one \texttt{proof} environment of 17 lines). No mathematics is written, repaired or re-derived: the statement and the proof are the article's own lines, in the article's own notation, and no retained line is substituted or re-flowed.  Retained uncounted as the source-local context the proof needs: lines 268--270; lines 339--341.  Nothing was omitted from any retained span, so the package carries no omission record.  No retained line contains a citation, so the wrapper carries no bibliography and no bibliography entry.  No retained line contains a figure environment, an image-inclusion command or drawing code, so no image is delivered and the package declares no asset dependency.  Editorial formatting only, in addition: an AMS \texttt{amsart} preamble that restates the article's own declarations and macros used by the retained lines, namely the environments \texttt{theorem} on separate counters, together with the standard packages those lines need.  The retained environments print in this document as Theorem 1 in order of appearance, whereas the article prints the same items with the numbering of its own full document.  The complete native source of the article as distributed in the arXiv e-print for this exact version is bundled unchanged as \texttt{sources/210001/source0.tex}.
\begin{equation}\label{form:ramanujan_1}
1 + 2 \sum_{n=1}^\infty \operatorname{sech}\left(\frac{K'(k)}{K(k)} \pi n\right) = \frac{2}{\pi} K(k), \quad \mathrm{with} \quad K(k) = \int_0^{\pi/2} \frac{\mathrm{d}\vartheta}{\sqrt{1-k^2\sin^2(\vartheta)}},
\end{equation}

\begin{equation}\label{eq:functional_F}
F(x) := \sum_{k=1}^\infty \frac{e^{-kx}}{1 + e^{-2kx}} =  \frac{\pi}{4x} - \frac{1}{4} + \frac{\pi}{x}F\left(\frac{\pi^2}{x}\right)
\end{equation}

\begin{theorem}
Let $x \in \mathbb{R}$. Then one has the functional equation
\begin{equation}
\sum_{k \in \mathbb{Z}} \operatorname{sech}\left(xk\right) = \frac{\pi}{x} \sum_{k \in \mathbb{Z}} \operatorname{sech}\left(\frac{\pi^2}{x} k \right)
\end{equation}
which is equivalent to (\ref{eq:functional_F}). For every $r \in \mathbb{N}$, this yields the probability mass function
$$
\left(q_{r,k} \right)_{k \in \mathbb{Z}} = \left(\frac{\pi}{\sqrt{r} \cdot 2K(k_r)} \operatorname{sech}\left( \frac{\pi}{\sqrt{r}} k \right) \right)_{k \in \mathbb{Z}}.
$$
for the discrete bilateral hyperbolic secant distribution with scaling factor $1/\sqrt{r}$.
\end{theorem}

\begin{proof}
The $-2i \pi$-Fourier transform $\mathcal{F}$ of the hyperbolic secant is well-known, viz.
$$
\mathcal{F}[u \mapsto \operatorname{sech}(u)](t) = \pi \operatorname{sech} \left(\frac{\pi}{2}t\right)
$$
Hence, by the Poisson summation formula, for every $x \in \mathbb{R}$ we get
\begin{align*}
\sum_{k \in \mathbb{Z}} \operatorname{sech}(xk) &=\ \sum_{k \in \mathbb{Z}} \mathcal{F}\left[y \mapsto \operatorname{sech}(xy)\right](\eta)\big|_{\eta = k}\\
&=\ \sum_{k\in \mathbb{Z}} \int_\mathbb{R} \operatorname{sech}(u) e^{-i(2\pi/x) \cdot \eta} \frac{\mathrm{d}u}{x}\big|_{\eta = k}\\
&=\ \frac{\pi}{|x|} \sum_{k \in \mathbb{Z}} \operatorname{sech}\left(\frac{\pi^2}{x}\eta\right) \big|_{\eta = k},
\end{align*}
which results in Equation (\ref{eq:functional_F}). Now, we substitute $x = \sqrt{r} \cdot \pi$. For every $r \in \mathbb{N}$, Ramanujan's formula \ref{form:ramanujan_1} yields
$$
\sum_{k\in \mathbb{Z}} \operatorname{sech}\left(\frac{\pi}{\sqrt{r}} k \right) = \sqrt{r} \left(\sum_{k \in \mathbb{Z}} \operatorname{sech}\left(\sqrt{r} \cdot \pi k \right)\right) = \frac{2\sqrt{r}}{\pi}K\left(k_r\right),
$$
as claimed.
\end{proof}

\end{document}
